\documentclass[conference, 10pt, twocolumn]{IEEEtran}

% ── Packages ──────────────────────────────────────────────────────────
\usepackage[utf8]{inputenc}
\usepackage[T1]{fontenc}
\usepackage{amsmath,amssymb}
\usepackage{graphicx}
\usepackage{booktabs}
\usepackage{multirow}
\usepackage{hyperref}
\usepackage{xcolor}
\usepackage{enumitem}
\usepackage{microtype}
\usepackage{tabularx}
\usepackage{caption}
\usepackage[numbers]{natbib}
\usepackage{float}

\hypersetup{
    colorlinks=true,
    linkcolor=blue!60!black,
    citecolor=blue!60!black,
    urlcolor=blue!60!black,
}

% ── Title ─────────────────────────────────────────────────────────────
\title{Multi-Task Financial NLP: Stance and Sentiment Classification\\
       of Central Bank Communications and News Headlines}

\author{
    \IEEEauthorblockN{COMP6713 Natural Language Processing -- Group Project}
    \IEEEauthorblockA{School of Computer Science and Engineering\\
    University of New South Wales\\
    Sydney, Australia}
}

\begin{document}
\maketitle

% ══════════════════════════════════════════════════════════════════════
%  ABSTRACT
% ══════════════════════════════════════════════════════════════════════
\begin{abstract}
Financial markets react rapidly to both central bank policy signals and
news sentiment.  In this project we tackle two closely related
text-classification tasks: (1) \emph{stance classification}---labelling
sentences from Federal Open Market Committee (FOMC) communications as
hawkish, dovish, or neutral---and (2) \emph{sentiment classification}---labelling
financial news headlines as positive, negative, or neutral.  We
systematically compare traditional machine-learning baselines (TF-IDF,
GloVe, lexicon rule-based), zero-shot and few-shot pre-trained
transformers, single-task fine-tuning, and a multi-task learning
architecture.  Our best system, a \textbf{multi-task FinBERT} model with
a shared encoder and dual task-specific heads, achieves 98.45\% accuracy
/ 0.9772 macro-F1 on sentiment and 67.74\% accuracy / 0.6684 macro-F1
on stance---outperforming all single-task alternatives.  We provide
quantitative evaluation, per-class analysis, error diagnostics, and a
Gradio demonstration interface.
\end{abstract}

% ══════════════════════════════════════════════════════════════════════
\section{Introduction}
\label{sec:intro}
% ══════════════════════════════════════════════════════════════════════

Central banks such as the U.S.\ Federal Reserve wield enormous influence
over global financial markets---not only through their policy actions,
but increasingly through the language they use to communicate those
actions.  A single hawkish or dovish sentence in an FOMC statement can
shift bond yields within seconds.  Likewise, the sentiment embedded in
financial news headlines shapes investor behaviour and drives short-term
price movements.

Manually analysing the volume and subtlety of these texts is
infeasible in real time.  Automated classification therefore promises
significant practical value: traders can react faster, risk managers can
monitor tone shifts continuously, and policymakers can study how their
communications are perceived.

In this report we address two complementary classification tasks:

\begin{enumerate}[leftmargin=*,nosep]
    \item \textbf{Stance classification} of FOMC sentences into
          \emph{hawkish} (favouring tighter monetary policy),
          \emph{dovish} (favouring looser policy), or \emph{neutral}.
    \item \textbf{Sentiment classification} of financial news headlines
          into \emph{positive}, \emph{negative}, or \emph{neutral}.
\end{enumerate}

Both are three-class text-classification problems, but they differ
substantially in difficulty: financial sentiment is relatively well
aligned with the prior knowledge of pre-trained language models, whereas
monetary-policy stance requires domain-specific understanding of subtle
economic signals.

Our contributions are as follows:
\begin{itemize}[leftmargin=*,nosep]
    \item A rigorous comparison of 14 model configurations spanning
          lexicon rules, classical ML, zero-shot/few-shot transformers,
          and fine-tuned transformers.
    \item A multi-task FinBERT architecture that jointly learns both
          tasks and achieves state-of-the-art results on both datasets.
    \item Detailed error analysis identifying the linguistic sources of
          remaining classification errors.
    \item A deployable Gradio web interface and command-line tool for
          real-time inference.
\end{itemize}

% ══════════════════════════════════════════════════════════════════════
\section{Related Work}
\label{sec:related}
% ══════════════════════════════════════════════════════════════════════

\textbf{Financial sentiment analysis} has progressed from dictionary-based
methods \cite{loughran2011liability} to transformer-based models.
Malo et al.\ introduced the Financial PhraseBank, demonstrating that
annotator agreement strongly affects benchmark quality.  Pre-trained
language models such as BERT \cite{devlin2019bert} have been adapted
to financial text through continued pre-training on financial corpora,
producing models such as FinBERT \cite{araci2019finbert,yang2020finbert},
which achieve strong zero-shot sentiment performance.

\textbf{Monetary-policy stance detection} is a newer task.
Shah et al.\ \cite{shah2023trillion} released the FOMC
Hawkish-Dovish dataset and showed that domain-adapted transformers
outperform general-purpose models.  The key challenge is that hawkish
and dovish language is often subtle and context-dependent---phrases such
as ``sustained progress toward our goals'' can be neutral or dovish
depending on context.

\textbf{Multi-task learning} (MTL) has been shown to improve
generalisation across related NLP tasks by encouraging shared
representations \cite{liu2019mt-dnn}.  To our knowledge, no prior
work jointly trains stance and sentiment in the financial domain.

% ══════════════════════════════════════════════════════════════════════
\section{Datasets}
\label{sec:data}
% ══════════════════════════════════════════════════════════════════════

\subsection{FOMC Hawkish-Dovish Dataset}
Released by the Georgia Tech Financial Technology Lab, this corpus
contains \textbf{2,480 sentences} from FOMC meeting minutes, speeches,
and press conferences, each labelled as \emph{hawkish}, \emph{dovish},
or \emph{neutral}.  The label distribution is imbalanced: approximately
50\% of sentences are neutral, with the remainder split unevenly between
hawkish and dovish.  This imbalance motivates our use of weighted
cross-entropy during training.

\subsection{Financial PhraseBank}
Introduced by Malo et al., this dataset contains \textbf{2,264 sentences}
drawn from financial news, labelled as \emph{positive}, \emph{negative},
or \emph{neutral}.  We use the subset where 100\% of annotators agreed,
yielding high-quality, unambiguous labels.

\subsection{Loughran-McDonald Financial Sentiment Dictionary}
This domain-specific lexicon categorises financial terms into sentiment
categories (positive, negative, uncertain, litigious, etc.).  We use it
in two ways: (a) as the basis for a rule-based baseline and (b) as
auxiliary features concatenated with TF-IDF representations.

\subsection{Data Splits}
Both datasets are split \textbf{70\% / 10\% / 20\%} for
training, validation, and testing, respectively.  We employ
\emph{stratified} splitting to preserve class proportions across
partitions, which is critical given the class imbalance in the FOMC
data.

% ══════════════════════════════════════════════════════════════════════
\section{Methodology}
\label{sec:method}
% ══════════════════════════════════════════════════════════════════════

We adopt a progressive modelling strategy, advancing from simple
baselines to increasingly expressive architectures.  This allows us to
quantify the value added by each modelling decision.

\subsection{Baseline Models}

\subsubsection{Lexicon Rule-Based (LM Lexicon)}
Each sentence is scored using the Loughran-McDonald dictionary: the
number of positive and negative terms is counted, and the sentence is
classified by the dominant count (ties default to neutral).  This
approach requires no training data and serves as a lower bound on
performance.

\subsubsection{TF-IDF + Logistic Regression}
We compute term frequency--inverse document frequency (TF-IDF) vectors
with unigram and bigram features and train an $\ell_2$-regularised
logistic regression classifier.  TF-IDF captures the discriminative
power of individual words and word pairs, while logistic regression
provides a well-calibrated linear decision boundary.  This serves as
our primary baseline.

\subsubsection{Alternative Baselines}
To understand the sensitivity of results to modelling choices, we
evaluate several variants:
\begin{itemize}[leftmargin=*,nosep]
    \item \textbf{TF-IDF + SVM}: Replaces logistic regression with a
          support vector machine using an RBF kernel, which can capture
          non-linear boundaries in the TF-IDF feature space.
    \item \textbf{TF-IDF (trigrams) + LR}: Extends the TF-IDF
          vocabulary to include trigrams, capturing longer phrases at
          the cost of higher dimensionality.
    \item \textbf{TF-IDF + LM Lexicon features}: Augments TF-IDF
          vectors with counts from the Loughran-McDonald dictionary,
          injecting domain knowledge into the feature space.
    \item \textbf{GloVe + LR / SVM}: Replaces TF-IDF with 300-dimensional
          GloVe embeddings averaged over tokens, providing
          semantically richer but lower-dimensional representations.
\end{itemize}

\subsection{Pre-Trained Transformer Evaluation}

\subsubsection{Zero-Shot FinBERT}
FinBERT is a BERT model further pre-trained on financial text.  It ships
with a native sentiment-classification head, which we evaluate directly
on the Financial PhraseBank test set without any additional training.
For the stance task, we use FinBERT's sentiment labels as a proxy
(positive $\to$ dovish, negative $\to$ hawkish, neutral $\to$ neutral),
an imperfect but informative mapping.

\subsubsection{Few-Shot Classification ($k{=}16$)}
We sample $k{=}16$ labelled examples per class and fine-tune three models:
FinBERT, BERT-base, and RoBERTa-base.  This setting tests how
efficiently each model can learn from minimal supervision and reveals
the advantage of domain-specific pre-training.

\subsection{Fine-Tuned Models}

\subsubsection{FinBERT Fine-Tuned (Single-Task)}
We fine-tune FinBERT separately on each dataset for 5 epochs, using
the AdamW optimiser with a learning rate of $2 \times 10^{-5}$ and
linear warm-up.  For the FOMC dataset, we apply weighted cross-entropy
loss to counteract class imbalance, assigning higher weights to
under-represented classes.

\subsubsection{BERT-base with LLRD and Gradual Unfreezing}
To investigate whether a general-purpose model can match domain-specific
pre-training, we fine-tune BERT-base with two advanced training
strategies:
\begin{itemize}[leftmargin=*,nosep]
    \item \textbf{Layer-wise Learning Rate Decay (LLRD)}: Lower layers
          (encoding general linguistic knowledge) receive smaller
          learning rates, while upper layers (encoding task-specific
          features) receive larger ones.  We use a decay factor of 0.9
          per layer.
    \item \textbf{Gradual unfreezing}: Layers are progressively
          unfrozen from top to bottom during training, preventing
          catastrophic forgetting of pre-trained representations.
    \item \textbf{Label smoothing} ($\alpha{=}0.1$): Softens the target
          distribution, improving calibration and reducing overfitting.
\end{itemize}
Training runs for 10 epochs with early stopping on validation macro-F1.

\subsection{Multi-Task FinBERT}
\label{sec:multitask}

Our best-performing architecture is a multi-task FinBERT model that
jointly learns stance and sentiment classification.  The design
rationale is that both tasks involve interpreting tone in financial
language, so a shared encoder should benefit from complementary
training signals.

\textbf{Architecture.}  The model consists of a shared FinBERT encoder
followed by two independent classification heads---one for stance and
one for sentiment.  Each head is a two-layer feedforward network with
dropout, taking the \texttt{[CLS]} token representation as input.

\textbf{Training procedure.}  We use alternating mini-batches: each
training step samples a batch from one task, computes the task-specific
loss, and updates both the shared encoder and the relevant head.
Batches from the two tasks are interleaved, ensuring balanced gradient
contributions.  We train for 8 epochs with AdamW ($\text{lr} = 2
\times 10^{-5}$, weight decay $= 0.01$).

\textbf{Loss function.}  Weighted cross-entropy is applied to the FOMC
task to address class imbalance.  The sentiment task uses standard
cross-entropy.  The total loss at each step is simply the loss of the
active task (no explicit task-weighting hyperparameter is needed because
alternating batches implicitly balance the two objectives).

% ══════════════════════════════════════════════════════════════════════
\section{Experiments}
\label{sec:experiments}
% ══════════════════════════════════════════════════════════════════════

\subsection{Evaluation Metrics}
We evaluate all models using:
\begin{itemize}[leftmargin=*,nosep]
    \item \textbf{Accuracy}: Overall fraction of correctly classified
          examples.
    \item \textbf{Macro-averaged F1}: The unweighted mean of per-class
          F1 scores, treating all classes equally regardless of size.
          This is our primary metric because it penalises poor
          performance on minority classes.
    \item \textbf{Per-class precision, recall, and F1}: Enables
          fine-grained diagnosis of which classes are confused.
    \item \textbf{Confusion matrices}: Visualise the pattern of
          misclassifications.
\end{itemize}

\subsection{Implementation Details}
All transformer models use the HuggingFace \texttt{transformers}
library.  Classical ML baselines use \texttt{scikit-learn}.  GloVe
embeddings are 300-dimensional vectors pre-trained on Common Crawl.
Experiments are conducted on a single NVIDIA GPU.  Random seeds are
fixed for reproducibility.

% ══════════════════════════════════════════════════════════════════════
\section{Results}
\label{sec:results}
% ══════════════════════════════════════════════════════════════════════

Tables~\ref{tab:sentiment} and~\ref{tab:stance} present the full
results for sentiment and stance classification, respectively.

% ── Sentiment Table ───────────────────────────────────────────────────
\begin{table}[t]
\centering
\caption{Sentiment classification results on Financial PhraseBank
(100\% annotator agreement).}
\label{tab:sentiment}
\small
\begin{tabular}{@{}lcc@{}}
\toprule
\textbf{Model} & \textbf{Acc.} & \textbf{M-F1} \\
\midrule
\multicolumn{3}{@{}l}{\textit{Traditional baselines}} \\
LM Lexicon (rule-based)     & 0.6932 & 0.5315 \\
TF-IDF + LR                 & 0.8720 & 0.8232 \\
TF-IDF + SVM                & 0.8940 & 0.8534 \\
TF-IDF (trigrams) + LR      & 0.8786 & 0.8310 \\
TF-IDF + LM Lexicon         & 0.8543 & 0.8050 \\
GloVe + LR                  & 0.7903 & 0.7218 \\
GloVe + SVM                 & 0.8035 & 0.7249 \\
\midrule
\multicolumn{3}{@{}l}{\textit{Pre-trained (no fine-tuning on full data)}} \\
FinBERT (zero-shot native)   & 0.9735 & 0.9650 \\
FinBERT (few-shot $k{=}16$)  & 0.9801 & 0.9690 \\
BERT-base (few-shot $k{=}16$)& 0.7461 & 0.6599 \\
RoBERTa (few-shot $k{=}16$)  & 0.7572 & 0.6439 \\
\midrule
\multicolumn{3}{@{}l}{\textit{Fine-tuned}} \\
FinBERT (fine-tuned)         & 0.9669 & 0.9467 \\
BERT-base LLRD               & 0.9691 & 0.9533 \\
\midrule
\multicolumn{3}{@{}l}{\textit{Multi-task}} \\
\textbf{Multi-task FinBERT}  & \textbf{0.9845} & \textbf{0.9772} \\
\bottomrule
\end{tabular}
\end{table}

% ── Stance Table ──────────────────────────────────────────────────────
\begin{table}[t]
\centering
\caption{Stance classification results on FOMC Hawkish-Dovish dataset.}
\label{tab:stance}
\small
\begin{tabular}{@{}lcc@{}}
\toprule
\textbf{Model} & \textbf{Acc.} & \textbf{M-F1} \\
\midrule
\multicolumn{3}{@{}l}{\textit{Traditional baselines}} \\
LM Lexicon (rule-based)     & 0.4153 & 0.3885 \\
TF-IDF + LR                 & 0.6089 & 0.5873 \\
TF-IDF + SVM                & 0.6331 & 0.6061 \\
TF-IDF (trigrams) + LR      & 0.6109 & 0.5914 \\
TF-IDF + LM Lexicon         & 0.6109 & 0.5863 \\
GloVe + LR                  & 0.5161 & 0.5040 \\
GloVe + SVM                 & 0.5544 & 0.5225 \\
\midrule
\multicolumn{3}{@{}l}{\textit{Pre-trained (no fine-tuning on full data)}} \\
FinBERT (zero-shot proxy)    & 0.4980 & 0.4874 \\
FinBERT (few-shot $k{=}16$)  & 0.4859 & 0.4552 \\
BERT-base (few-shot $k{=}16$)& 0.3790 & 0.3694 \\
RoBERTa (few-shot $k{=}16$)  & 0.3589 & 0.3489 \\
\midrule
\multicolumn{3}{@{}l}{\textit{Fine-tuned}} \\
FinBERT (fine-tuned)         & 0.6129 & 0.5988 \\
BERT-base LLRD               & 0.6512 & 0.6371 \\
\midrule
\multicolumn{3}{@{}l}{\textit{Multi-task}} \\
\textbf{Multi-task FinBERT}  & \textbf{0.6774} & \textbf{0.6684} \\
\bottomrule
\end{tabular}
\end{table}

% ══════════════════════════════════════════════════════════════════════
\section{Analysis}
\label{sec:analysis}
% ══════════════════════════════════════════════════════════════════════

\subsection{Sentiment vs.\ Stance: A Tale of Two Tasks}

The most striking pattern in our results is the large gap between the
two tasks.  The best sentiment model achieves 98.45\% accuracy, whereas
the best stance model reaches only 67.74\%.  This gap persists across
\emph{every} model family and reflects a fundamental difference in task
difficulty.

Financial sentiment is closely aligned with general language semantics:
words like ``profit'', ``growth'', and ``loss'' carry clear sentiment
signals.  Monetary-policy stance, by contrast, involves interpreting
domain-specific nuance.  A sentence such as ``\emph{The Committee
judges that the risks to the economic outlook are roughly balanced}''
is neutral, but a subtle rewording---``\emph{The Committee judges that
the risks are weighted to the upside}''---becomes hawkish.  This kind
of fine-grained distinction is difficult even for human annotators.

\subsection{The Value of Domain-Specific Pre-Training}

FinBERT's performance in zero-shot and few-shot settings is remarkable.
With \emph{zero} task-specific training, FinBERT achieves 97.35\%
accuracy on sentiment, outperforming every traditional baseline.  In
the few-shot setting ($k{=}16$), FinBERT achieves 98.01\% on sentiment
versus 74.61\% for BERT-base and 75.72\% for RoBERTa---a gap of over
20 percentage points.

This demonstrates that domain-specific pre-training on financial corpora
equips the model with strong priors about financial language.  General-purpose
models, despite their larger training corpora, lack the specialised
vocabulary and contextual understanding needed for financial text.

On the harder stance task, however, zero-shot and few-shot approaches
fall short.  FinBERT's zero-shot proxy achieves only 49.80\% accuracy
(near random for a three-class problem), and few-shot performance is
similarly poor across all models.  This confirms that stance detection
fundamentally requires supervised fine-tuning.

\subsection{Multi-Task Learning Benefits}

The multi-task FinBERT achieves the best results on \emph{both} tasks
simultaneously:
\begin{itemize}[leftmargin=*,nosep]
    \item \textbf{Sentiment}: 98.45\% accuracy (+1.76 pp over
          single-task FinBERT).
    \item \textbf{Stance}: 67.74\% accuracy (+6.45 pp over single-task
          FinBERT).
\end{itemize}

The improvement is larger on stance, suggesting that sentiment
supervision provides useful inductive bias for the harder task.
Intuitively, understanding whether financial language is positive or
negative helps disambiguate hawkish and dovish signals.

\subsection{BERT-base LLRD: Closing the Gap Without Domain Pre-Training}

An interesting finding is that BERT-base with LLRD and gradual
unfreezing achieves 65.12\% accuracy on stance---surpassing
FinBERT's single-task fine-tuned result of 61.29\%.  This shows that
careful training strategies (layer-wise learning rate decay, gradual
unfreezing, label smoothing) can compensate for the lack of
domain-specific pre-training.  On sentiment, BERT-base LLRD matches
FinBERT fine-tuned at 96.91\%, confirming that the training regime
matters as much as the starting checkpoint for well-aligned tasks.

\subsection{Traditional Baselines}

Among traditional methods, TF-IDF + SVM is the strongest, achieving
89.40\% on sentiment and 63.31\% on stance.  The SVM's ability to find
non-linear boundaries in high-dimensional TF-IDF space gives it a
consistent edge over logistic regression.  Extending to trigrams provides
marginal improvement, while adding Loughran-McDonald lexicon features
actually \emph{hurts} TF-IDF performance---likely because the lexicon
introduces noise when its categories do not align perfectly with the
task labels.

GloVe-based models underperform TF-IDF, averaging the embeddings of all
tokens loses the discriminative power of rare but informative words that
TF-IDF naturally upweights.

\subsection{Error Analysis}
\label{sec:errors}

We conduct a qualitative analysis of misclassified examples from the
multi-task FinBERT model.

\subsubsection{Stance Errors: Neutral vs.\ Hawkish Confusion}
The most common error pattern on the FOMC dataset is confusion between
neutral and hawkish sentences.  Many FOMC communications use hedged
language that conveys a hawkish lean without explicit hawkish keywords.
For example:

\begin{quote}
\small
``\emph{Participants noted that inflation expectations remained well
anchored, although several observed that the recent data warranted
close monitoring.}''
\end{quote}

This sentence has a subtle hawkish undertone (``warranted close
monitoring''), but the model classifies it as neutral because the
majority of the sentence discusses stable expectations.  Such errors
are difficult to resolve without richer contextual modelling (e.g.,
document-level features or discourse structure).

\subsubsection{Sentiment Errors: Neutral vs.\ Positive Ambiguity}
On the Financial PhraseBank, the few errors tend to involve factual
statements with mildly positive implications:

\begin{quote}
\small
``\emph{The company reported revenue of \$3.2 billion, in line with
analyst estimates.}''
\end{quote}

Meeting estimates is arguably neutral, but some annotators may
interpret ``in line'' as mildly positive.  The model occasionally
misclassifies these borderline cases.

\subsection{Effect of Weighted Cross-Entropy}

Applying class-weighted loss to the FOMC task is critical.  Without
weighting, the model learns to default to the majority ``neutral''
class, achieving high accuracy but poor macro-F1.  With weighting,
the model is incentivised to correctly classify hawkish and dovish
sentences, improving macro-F1 significantly even if overall accuracy
dips marginally on the majority class.

% ══════════════════════════════════════════════════════════════════════
\section{Deployment}
\label{sec:deployment}
% ══════════════════════════════════════════════════════════════════════

To demonstrate practical utility, we provide two inference interfaces:

\begin{itemize}[leftmargin=*,nosep]
    \item \textbf{Command-Line Interface (CLI)}: Users can pass a
          sentence and receive stance and sentiment predictions with
          confidence scores.
    \item \textbf{Gradio Web Interface}: A browser-based demo that
          allows interactive input and displays predicted labels
          alongside model confidence for each class.
\end{itemize}

Both interfaces load the multi-task FinBERT checkpoint and perform
inference in real time on CPU or GPU.

% ══════════════════════════════════════════════════════════════════════
\section{Discussion}
\label{sec:discussion}
% ══════════════════════════════════════════════════════════════════════

\textbf{Why Multi-Task?}  The results confirm our hypothesis that
stance and sentiment share underlying structure.  Both tasks require
the model to interpret tone in financial language, and the shared
encoder benefits from seeing both types of supervision.  The
alternating-batch strategy is simple yet effective, avoiding the need
for task-weighting hyperparameters.

\textbf{Limitations.}  Our stance accuracy of 67.74\% leaves room for
improvement.  Potential directions include: (a) incorporating
document-level context (FOMC sentences are currently classified in
isolation); (b) using larger models (e.g., LLaMA or GPT-4) for
knowledge distillation; (c) data augmentation to address class
imbalance more aggressively; and (d) integrating temporal information,
since monetary-policy language evolves across meeting cycles.

\textbf{Broader Impact.}  Automated stance and sentiment tools could
amplify herd behaviour in financial markets if adopted uncritically.
We recommend that such systems be used as decision-support tools
rather than autonomous trading signals.

% ══════════════════════════════════════════════════════════════════════
\section{Conclusion}
\label{sec:conclusion}
% ══════════════════════════════════════════════════════════════════════

We presented a systematic study of financial text classification
spanning rule-based lexicons, classical ML, and transformer-based
models.  Our key findings are:

\begin{enumerate}[leftmargin=*,nosep]
    \item \textbf{Domain pre-training matters}: FinBERT dramatically
          outperforms general-purpose BERT and RoBERTa in low-data
          settings, especially for sentiment.
    \item \textbf{Stance requires fine-tuning}: Unlike sentiment,
          monetary-policy stance cannot be solved by zero-shot or
          few-shot transfer; supervised learning on labelled FOMC
          data is essential.
    \item \textbf{Multi-task learning helps}: Jointly training on
          stance and sentiment yields the best results on both tasks,
          with a larger improvement on the harder stance task.
    \item \textbf{Training strategies bridge the gap}: BERT-base with
          LLRD and gradual unfreezing achieves competitive stance
          performance without domain-specific pre-training.
\end{enumerate}

Our multi-task FinBERT model achieves 98.45\% accuracy on sentiment and
67.74\% on stance, representing the best results across all experiments.
We release a Gradio demo and CLI tool for interactive exploration.

% ══════════════════════════════════════════════════════════════════════
%  REFERENCES
% ══════════════════════════════════════════════════════════════════════

\bibliographystyle{plainnat}

\begin{thebibliography}{10}

\bibitem{loughran2011liability}
T.~Loughran and B.~McDonald,
``When is a liability not a liability? Textual analysis, dictionaries,
and 10-Ks,''
\emph{The Journal of Finance}, vol.~66, no.~1, pp.~35--65, 2011.

\bibitem{devlin2019bert}
J.~Devlin, M.-W. Chang, K.~Lee, and K.~Toutanova,
``BERT: Pre-training of deep bidirectional transformers for language
understanding,''
in \emph{Proc.\ NAACL-HLT}, 2019, pp.~4171--4186.

\bibitem{araci2019finbert}
D.~Araci,
``FinBERT: Financial sentiment analysis with pre-trained language
models,''
\emph{arXiv preprint arXiv:1908.10063}, 2019.

\bibitem{yang2020finbert}
Y.~Yang, M.~C.~S. Uy, and A.~Huang,
``FinBERT: A pretrained language model for financial communications,''
\emph{arXiv preprint arXiv:2006.08097}, 2020.

\bibitem{shah2023trillion}
A.~Shah, S.~Paturi, and S.~Chava,
``Trillion dollar words: A new financial dataset, task \& market
analysis,''
in \emph{Proc.\ ACL}, 2023.

\bibitem{liu2019mt-dnn}
X.~Liu, P.~He, W.~Chen, and J.~Gao,
``Multi-task deep neural networks for natural language understanding,''
in \emph{Proc.\ ACL}, 2019, pp.~4487--4496.

\bibitem{malo2014good}
P.~Malo, A.~Sinha, P.~Korhonen, J.~Wallenius, and P.~Takala,
``Good debt or bad debt: Detecting semantic orientations in economic
texts,''
\emph{Journal of the Association for Information Science and Technology},
vol.~65, no.~4, pp.~782--796, 2014.

\bibitem{howard2018ulmfit}
J.~Howard and S.~Ruder,
``Universal language model fine-tuning for text classification,''
in \emph{Proc.\ ACL}, 2018, pp.~328--339.

\bibitem{liu2019roberta}
Y.~Liu, M.~Ott, N.~Goyal, J.~Du, M.~Joshi, D.~Chen, O.~Levy,
M.~Lewis, L.~Zettlemoyer, and V.~Stoyanov,
``RoBERTa: A robustly optimized BERT pretraining approach,''
\emph{arXiv preprint arXiv:1907.11692}, 2019.

\bibitem{pennington2014glove}
J.~Pennington, R.~Socher, and C.~D.~Manning,
``GloVe: Global vectors for word representation,''
in \emph{Proc.\ EMNLP}, 2014, pp.~1532--1543.

\end{thebibliography}

\end{document}
